\begin{lemma}[Exponential palette cost of minimal infinite directions]
\label{lem:rh-direction-palette-bound}
Let $F=\Sigma(\Id)$ be the greatest finite-class ORIGINAL invariant
equivalence and let $E_1,\ldots,E_r$ list all minimal infinite
ORIGINAL invariant equivalences (with $r=0$ if $M$ is finite).
Here and in the next two corollaries we abbreviate
$\mathcal D\setminus F=\{a\in\mathcal D:a\not\subseteq F\}$
and
$\mathcal D\cap F=\{a\in\mathcal D:a\subseteq F\}$
as SETS OF ATOM NAMES, not literal set operations between
binary relations and names.
Then
\[
   2^r-1\le|\mathcal D\setminus F|,
   \qquad
   r\le\left\lfloor
       \log_2\bigl(|\mathcal D\setminus F|+1\bigr)
                         \right\rfloor.
\]
All the counted colours are distinct ACTUAL original diversity
orbitals outside $F$; $F$ is NOT assumed exact.
\end{lemma}
\begin{proof}
Let $\mathcal L$ be the finite distributive lattice of original
FULL RAW exact equivalences, by
\cref{lem:acf-exact-distributive}.  Let $K_{\mathrm{fin}}$
be the join in $\mathcal L$ of ALL finite-class exact
equivalences, including $\Id$.  This join has finite classes:
the product of commuting finite-class equivalences has finite
classes, and the finite exact lattice is closed under joins.
Thus $K_{\mathrm{fin}}\le F$, but equality need not hold.

Each $E_i$ is exact by \cref{prop:minimalcut}.  Put
$A_i=K_{\mathrm{fin}}\vee E_i$ and, for
$S\subseteq[r]$, put
$A_S=K_{\mathrm{fin}}\vee\bigvee_{i\in S}E_i$, with
$A_\varnothing=K_{\mathrm{fin}}$.
For distinct $i,j$, the intersection $E_i\cap E_j$
is a finite-class exact equivalence by minimality of the
infinite $E_i$, hence is contained in $K_{\mathrm{fin}}$.
Distributivity inside $\mathcal L$ gives
\[
 A_i\cap A_j=K_{\mathrm{fin}},\qquad
 A_S\cap A_T=A_{S\cap T}\quad(S,T\subseteq[r]).
\tag{PD}\label{eq:rh-direction-boolean}
\]
These equalities use only the exact sublattice, NOT the
possibly nondistributive full invariant-equivalence lattice.

For nonempty $S$ and $i\in S$, the $A_i$-classes are infinite,
the $K_{\mathrm{fin}}$-classes are finite, and
$A_i\cap A_{S\setminus\{i\}}=K_{\mathrm{fin}}$.
Thus a single $A_i$-class contains infinitely many
distinct $A_{S\setminus\{i\}}$-classes.  Therefore
\[
 [A_S:A_{S\setminus\{i\}}]=\infty.
\]
Likewise $[A_S:F\cap A_S]=\infty$, because the former
has infinite classes and the latter has finite classes.

Fix one real $x\in M$, put $G=\Aut(M)$, and for every
original invariant equivalence $H$ write
$G_H=\{g\in G:xH gx\}$.
This is a subgroup, and for $H\le K$ the standard
transitive $G$-set correspondence gives
$[G_K:G_H]=[K:H]$.  Hence, INSIDE $G_{A_S}$,
each of the finitely many subgroups
$G_{A_{S\setminus\{i\}}}$ ($i\in S$) and
$G_{F\cap A_S}$ has infinite index.  By B.~H.~Neumann's
finite-coset-cover theorem~\cite{NeumannCover}, also used in
\cref{lem:finite-action-avoidance}, their union cannot
cover $G_{A_S}$.  Choose $g_S$ outside their union.
The old colour $t_S=\col(x,g_Sx)$ lies in $A_S$,
but not in $F$ or in any $A_{S\setminus\{i\}}$.

If $S,T$ are distinct nonempty subsets, say
$i\in S\setminus T$, a putative equality $t_S=t_T$
would place this single old atom in
$A_S\cap A_T=A_{S\cap T}\subseteq A_{S\setminus\{i\}}$,
contrary to its choice.  When $S\subsetneq T$, reverse
the roles.  Thus the $2^r-1$ nonempty subsets yield
pairwise distinct original diversity colours outside $F$.
\end{proof}

\begin{corollary}[Bounded character of real algebraic closure]
\label{cor:rh-acl-bounded-character}
Let $E_1,\ldots,E_r$ list all inclusion-minimal ORIGINAL invariant
equivalences with infinite real classes.  Then
$r\le\lfloor\log_2(|\mathcal D\setminus F|+1)\rfloor\le|\mathcal D|$
by \cref{lem:rh-direction-palette-bound}, and
for every finite real $A\subseteq M$,
\[
   \operatorname{acl}_M(A)=
   \bigcup_{\substack{B\subseteq A\\|B|\le r}}
                 \operatorname{acl}_M(B).
\]
For $x\notin A$, this follows from the exact colour criterion
\[
 x\in\operatorname{acl}_M(A)
 \quad\Longleftrightarrow\quad
 \text{for every }i\le r\text{ some }a\in A
 \text{ satisfies }E_i\nsubseteq
       E_{\Gamma(\col(x,a))}.
\]
Here $\Gamma(p)=\langle\{p\}\rangle$ is the ORIGINAL generated
equivalence support, not an exact quotient kernel.  If $M$ is
finite then $r=0$ and $\operatorname{acl}_M(\varnothing)=M$.
\end{corollary}
\begin{proof}
For $x\notin A$, set
$J=\bigcap_{a\in A}\Gamma(\col(x,a))$, with empty intersection
the universal support.  Since
$xE_{\Gamma(\col(x,a))}a$ for each $a$, its actual $E_J$-class is
\[
 [x]_{E_J}=
       \bigcap_{a\in A}[a]_{E_{\Gamma(\col(x,a))}}.
\]
If $x$ is algebraic over nonempty $A$, this class is finite by
Theorem~\ref{thm:rh-hull}; the empty-base case follows from
transitivity.  Conversely, if the class is finite, the orbit of
$x$ under $\Aut(M/A)$ is contained in it: each automorphism
fixing $A$ preserves the literal old spoke colours to every
$a\in A$.  Thus $x$ is algebraic over $A$.
Because the invariant-equivalence lattice is finite, $J$
has infinite classes exactly when it contains some minimal
infinite $E_i$.  Negating
$E_i\le J=\bigcap_a\Gamma(\col(x,a))$ gives the displayed
colour criterion.

For each $x\in\operatorname{acl}_M(A)\setminus A$, choose one
witness $a_i\in A$ excluding each $E_i$, and take the set $B$
of these witnesses.  Then $|B|\le r$ and the same criterion
gives $x\in\operatorname{acl}_M(B)$.  If $x\in A$ and $M$
is infinite, $r\ge1$ and $B=\{x\}$ suffices.  If $M$ is
finite all real points are algebraic over the empty set.
Monotonicity proves the reverse inclusion.

Finally let $F=\Sigma(\Id)$ be the greatest finite-class
invariant equivalence.  Each $E_i$ contains some old
diversity atom $t_i$ outside $F$, for otherwise it would
have finite classes.  For $i\ne j$ the proper intersection
$E_i\cap E_j$ has finite classes by minimality, hence lies
below $F$.  Thus $t_i\ne t_j$, giving the injective bound
$r\le|\mathcal D\setminus F|\le|\mathcal D|$.
For a singleton parameter $\{a\}$, the already established
colour criterion says
$x\in\operatorname{acl}_M(a)$ if and only if
$\Gamma(\col(x,a))\subseteq F$ for $x\ne a$,
which is equivalent to $xFa$.
Hence $\operatorname{acl}_M(a)=F[a]$.
If $r=1$, the bounded-character formula consequently yields
$\operatorname{acl}_M(A)=F[A]$ for every finite real $A$.
This alternative short derivation does NOT prove that naming
quotient owner imaginaries is harmless.
\end{proof}

\begin{corollary}[Polynomial growth and finite-packet bound for real algebraic closure]
\label{cor:rh-acl-polynomial-growth}
Assume $M$ is infinite and let $F=\Sigma(\Id)$.
Write $f=|\mathcal D\cap F|$,
$d_\infty=|\mathcal D\setminus F|$ and
$c=|[x]_{E_F}|$, the constant finite size of any
original $F$-class.  Then $c\le 2^f$, and for
every finite real $A$ of size $n$,
\[
 |\operatorname{acl}_M(A)|
 \le c\left(n+
      \sum_{j=1}^{\min\{r,n\}}
           \binom{n}{j}d_\infty^j\right),
 \qquad
 r\le\left\lfloor\log_2(d_\infty+1)\right\rfloor.
\]
In particular $|\operatorname{acl}_M(A)|=O_M(n^r)$.
Moreover $\operatorname{acl}_M(A)$ is a union of COMPLETE
original $F$-classes, and the number of those classes is at most
$n+\sum_{j=1}^{\min\{r,n\}}\binom{n}{j}d_\infty^j$.
If $F=\Id$, then
$\operatorname{acl}_M(A)=\operatorname{dcl}_M(A)$
for every finite REAL $A$, although this closure may
strictly contain $A$.
No analogous assertion for arbitrary imaginary sorts is made.
\end{corollary}
\begin{proof}
The invariant $F$-classes are finite and transitive;
\cref{cor:rh-class-2power} gives $c\le 2^f$.
By \cref{thm:rh-globalise}, every finite realised
complete one-point fibre over any finite real parameter
set is ONE class of an original global finite-class
equivalence $K\le F$, and so has size at most $c$.
The finite $F$-closure $F[A]$ contains at most $cn$ points.

For every $x\in\operatorname{acl}_M(A)\setminus F[A]$,
\cref{cor:rh-acl-bounded-character} gives a nonempty
$B\subseteq A$ with $|B|\le r$ and
$x\in\operatorname{acl}_M(B)$.
Since $x\notin F[A]$, every old spoke colour
$\col(x,b)$ ($b\in B$) lies OUTSIDE $F$.
For fixed $B$ of size $j$ there are at most
$d_\infty^j$ such complete original spoke-colour
tuples.  Each realised finite type fibre has
size at most $c$ by globalisation, and distinct
fibre types are disjoint.  Thus at most
$c\,d_\infty^j$ such points are algebraic over
$B$ outside $F[A]$.  Summing over the
$\binom nj$ choices of $B$ and adding $F[A]$
gives the displayed inequality, without presuming
disjointness between closures from different bases.
The sharp logarithmic upper bound on $r$ is
\cref{lem:rh-direction-palette-bound}.
The empty-$A$ case holds by infinite transitivity.
If $x\in\operatorname{acl}_M(A)$ and $xFy$, then
$y\in\operatorname{acl}_M(x)$, because the original $F$-class
of $x$ is finite and definable over $x$.
Transitivity of algebraicity gives
$y\in\operatorname{acl}_M(A)$.
Thus $\operatorname{acl}_M(A)$ is $F$-saturated;
each of its $F$-classes has size $c$, and dividing
the preceding cardinal bound by $c$ gives the packet count.

If $F=\Id$, globalisation makes every finite
one-point type fibre a singleton.  Its complete
finite quantifier-free spoke formula over $A$
then defines its unique realisation in $M$.
Thus every algebraic real point over $A$ is
definable; the converse is always true.
\end{proof}

\begin{example}[Sharp closure exponent in coordinate sources]
\label{ex:rh-acl-polynomial-sharp}
For each $r\ge2$ let $M_r=\mathbb N^r$ and name the symmetric
diversity atom of distinct $x,y$ by the nonempty difference
support $\{i:x_i\ne y_i\}\subseteq[r]$.  Its automorphism group
is $\prod_{i=1}^r\operatorname{Sym}(\mathbb N)$: named atoms
preserve each coordinate equality, while independent coordinate
permutations preserve all atoms.  Thus it is homogeneous.
Every triangle-legal complete diagram induces one equivalence
relation of coordinate equality for each $i$; these can be
realised independently on $\mathbb N$, and their meet is real
identity because every specified diversity support is nonempty.
Hence the age is triangle-defined and has AP.

It also has semantic homomorphic CCO: an inconsistent
coordinate equality/disequality specification has a specified
inequality edge whose endpoints are connected by a path
of equality edges in that coordinate.  A shortest such path
together with the inequality edge is a specified simple
cycle without a homomorphism.  If no such obstruction exists
in any coordinate, assign different values to the finitely many
equality components independently, giving a homomorphism of
the entire partial diagram.

Its invariant equivalences are exactly those allowing changes
within a fixed set of coordinates: if an equivalence contains
an atom with difference support $S$, a two-step path of that
same atom (using a fresh third value in coordinate $i$)
forces each singleton difference support $\{i\}$ for $i\in S$;
single-coordinate steps then generate the whole corresponding
coordinate-change equivalence.  Thus there are precisely $r$
minimal infinite equivalences.  For
$A_n=\{(k,\ldots,k):1\le k\le n\}$, the pointwise stabiliser fixes
every tuple in $[n]^r$ and has infinite orbit on any tuple
with a coordinate outside $[n]$.  Therefore
\[
 \operatorname{acl}_{M_r}(A_n)
 =\operatorname{dcl}_{M_r}(A_n)=[n]^r,\qquad
 |\operatorname{acl}_{M_r}(A_n)|=n^r.
\]
The polynomial exponent $r$ is therefore optimal even
for definable closure.
\end{example}

